\section{Context Fluency: An Emerging Skill}

Effective agentic coding requires a skill set distinct from both
traditional programming and prompt engineering. We propose
\emph{context fluency} to describe this capability: the ability to
create rich, structured context that AI agents can act on. Context
fluency captures domain expertise, value judgments, and design intent
in machine-legible form---not as one-shot prompts, but as persistent
informational environments. Prompt engineering optimizes individual
instructions; context fluency is \emph{upstream} of prompting,
concerned with the informational architecture surrounding agent
execution~\cite{anthropic2025context, gartner2025context}. Where MEP
describes a sequence of phases producing artifacts, context fluency
describes the practitioner skill that makes the process effective.

We identify four components of context fluency:
\begin{description}
  \item[Decomposition:] Breaking problems into discrete, parallelizable
    tasks with clear boundaries, enabling concurrent agent execution.
  \item[Specification:] Describing not just \emph{what} to build but
    \emph{why}, so agents can make aligned micro-decisions without human
    intervention.
  \item[Constraint definition:] Knowing what to exclude, simplify, or
    defer---scope management as a first-class concern.
  \item[Domain encoding:] Externalizing tacit
    knowledge~\cite{polanyi1966tacit} that agents cannot generate on
    their own---directly instantiating what Nonaka and
    Takeuchi~\cite{nonaka1995knowledge} term \emph{externalization},
    converting tacit knowledge into explicit, communicable form.
    Twenty minutes of dictated pedagogical intuitions
    produced context that substantially reduced iteration on the
    domain-specific tutor component.
\end{description}

These components map to established frameworks. Backward
design~\cite{wiggins1998understanding} contributes outcome-driven
preparation; Polanyi's tacit knowledge~\cite{polanyi1958personal}
frames the externalization challenge---``we know more than we can
tell,'' and context fluency is the discipline of telling it anyway.
Context fluency also resembles pedagogical scaffolding: structuring
an environment so the learner (or agent) can act independently.

If context fluency is a distinct skill, the implications extend
beyond individual practice. For developer education, curricula should
cultivate specification, decomposition, and domain-encoding alongside
programming competencies. For hiring, practitioners with strong
domain knowledge and pedagogical instincts may be disproportionately
effective in agentic workflows~\cite{dakhel2025humanai}. For tool
design, preparation aids---task systems~\cite{yegge2025beads,
emanuel2025beadsrust}, specification frameworks~\cite{github2025speckit},
and context engineering platforms~\cite{anthropic2025context}---are
essential infrastructure, not optional enhancements.
